\begin{ledgroupsized}[r]{120mm}
\footnotesize
\pstart
\noindent\textbf{Überlieferung:}
\pend
\end{ledgroupsized}
\begin{ledgroupsized}[r]{114mm}
\footnotesize
\pstart \parindent -6mm
\makebox[6mm][l]{\textit{L}}%
Konzept: LBr~724 Bl.~16\textendash16a.
Ein Bogen~4\textsuperscript{o}, Ränder beschnitten.
Zwei Seiten auf Bl.~16a.
Textfolge: 16a~r\textsuperscript{o} bis auf den Textbereich \glqq NB\grqq\ in der oberen Ecke, dann Bl.~16a~v\textsuperscript{o}, zuletzt der Bereich \glqq NB\grqq\ auf Bl.~16a~r\textsuperscript{o}.
Bl.~16 enthält Leibnizens Abschrift eines \cite{02010}Briefes von 
\indexN{\textls{Vanni}, Giovanni Francesco 1638\textendash1709}G.\,F.~Vanni an 
\indexN{\textls{Magliabechi} (Magliabecchi, Magliabecki, Magliabekius, Magliabequius), Antonio 1633\textendash1714}A.~Magliabechi 
vom 5. Juni 1685 (\textit{LSB} III, 4 N.\;97), sowie eine Nebenrechnung zum vorliegenden Stück (siehe \refnote{FnA:LBr_724_16a_Nebenrechnung}).
\pend
\end{ledgroupsized}
\pstart
\noindent
\href{\Scan{LBr_724_16ar}}%
{\lbrack16a~r\textsuperscript{o}\rbrack} 
\pend
\nobreak
\begin{edfigure}[0.5em]%
\includegraphics[scale=1]{%
gesamttex/edit_VIII-4/images/LBr_724_016a_d1_016ar.pdf}
\edcaption{\edtext{\lbrack\textit{Fig.~1}\rbrack}{\lemma{\lbrack\textit{Fig.~1}\rbrack}%
\Cfootnote[nonum]{Vorlage ist die Figur aus \indexN{\textls{Vanni}, Giovanni Francesco 1638\textendash1709}\textsc{G.\,F. Vanni},
\cite{01995}\textit{Specimen libri de momentis gravium}, %
\cite{01023}\textit{AE} November 1684, S.\;511\textendash513, hier S.\;512.
Ein gestrichener Entwurf zum Diagramm wird nicht wiedergegeben.}}}
\end{edfigure}
\pstart
\noindent
\edlabel{LBr_724_16a_4a}%
Ob duas causas 
\edtext{planum\indexS{planum declive} alterutrum}{\lemma{planum}\Bfootnote{%
\textit{(1)}~premi %
\textit{(a)}~intel %
\textit{(b)}~intelli %
\textit{(2)}~alterutrum~\textit{L}}} 
ut \textit{XC}, 
\edtext{a globo}{\lemma{a}\Bfootnote{%
\textit{(1)}~gravi %
\textit{(2)}~globus %
\textit{(3)}~globo~\textit{L}}} 
\textit{I} premi
\edtext{intelligitur. Prima est quod globus\indexS{globus in plano descendere tendens} \textit{I} descendere tendens}{\lemma{intelligitur}\Bfootnote{%
\textit{(1)}~nam %
\textit{(a)}~quia etiamsi abesset \bars{planum \textit{streicht Hrsg.}} \textit{ZC}, %
\textit{(aa)}~ipsumque %
\textit{(bb)}~et globus \textit{I} libere decurrere posset in plano 
\bars{\bars{declivi \textit{erg.}}
\textit{XC} \textit{streicht Hrsg.}}
tamen premeret planum \textit{XC}, eo ipso dum ab \bars{eo \textit{streicht Hrsg.}} sustentatur, ne %
\textit{(aaa)}~libere rec %
\textit{(bbb)}~recta descendere possit %
\textit{(2)}~. Prima est %
\textit{(a)}~quod sublato etiam plano \textit{ZC} tamen a globo %
\textit{(b)}~quod globus %
\textit{(aa)}~decurre %
\textit{(bb)}~\textit{I} %
\textit{(aaa)}~decurrere tendens %
\textit{(bbb)}~descendere incipiens \bars{in \textit{streicht Hrsg.}} aliquo plano declivi ut \textit{XC} %
\textit{(ccc)}~descendere tendens~\textit{L}}} 
in linea \textit{FC},
\edtext{momento\indexS{momentum} quod}{\lemma{momento}\Bfootnote{%
\textit{(1)}~ut %
\textit{(2)}~quod~\textit{L}}} 
sit ad totale,\indexS{momentum totale} ut \textit{XN} ad \textit{XC} 
\edlabel{LBr_724_16a_3a}%
\edtext{}{%
\xxref{LBr_724_16a_3a}{LBr_724_16a_3b}%
\lemma{(quemadmodum \lbrack...\rbrack\ demonstravimus)}%
\Cfootnote{%
Das vorliegende Konzept N.\;\ref{32310} enthält keinen Beweis des Theorems der Momente, aber ein solcher wird in N.\;\ref{32313} erbracht, das neben N.\;\ref{32310} eine Vorarbeit zur \textit{Demonstratio geometrica} (N.\;\ref{49409} in diesem Band), und zwar zu ihrem ersten Teil, darstellt.}}%
(quemadmodum 
\edtext{demonstravimus)\lbrack,\rbrack\ aget reliquo,}{\lemma{demonstravimus)}\Bfootnote{%
\textit{(1)}~agit reliquo, (quod erit ut $XC - XN$ %
\textit{(2)}~aget reliquo,~\textit{L}}}
(quod erit
\edtext{ad totale\indexS{momentum totale}}{\lemma{}\Bfootnote{ad totale \textit{erg.}~\textit{L}}}
ut $XC - XN$
\edtext{ad \textit{XC},)}{\lemma{}\Bfootnote{ad \textit{XC} \textit{erg.}~\textit{L}}}
in ipsum planum declive\indexS{planum declive}
\edtext{\textit{XFC}, a quo}{\lemma{\textit{XFC}}\Bfootnote{%
\textit{(1)}~. Altera est quod %
\textit{(2)}~, a quo~\textit{L}}}
sustentatur ne perpendiculariter descendere possit.
\pend 
\pstart
Secunda
\edtext{causa premendi planum\indexS{planum declive}}{\lemma{causa}\Bfootnote{%
\textit{(1)}~est, quod globus \textit{I} descendens in %
\textit{(a)}~\textit{HC} momento %
\textit{(b)}~plano %
\textit{(2)}~premendi %
\textit{(a)}~glo %
\textit{(b)}~planum~\textit{L}}}
\textit{XC} est, quod globus \textit{I} descendens in altero plano \textit{HC} momento quod sit ad totale ut \textit{ZO} ad \textit{ZC}
(quemadmodum demonstravimus)\edlabel{LBr_724_16a_3b}
premit hoc ipso momento planum \textit{XC}.
\pend 
\pstart
\edtext{Cumque eadem ratiocinatio institui possit}{\lemma{Cumque}\Bfootnote{%
\textit{(1)}~idem dici possit %
\textit{(2)}~eadem ratiocinatio institui possit~\textit{L}}} et de plano
\edtext{altero \textit{ZC} hinc ponendo}{\lemma{altero}\Bfootnote{%
\textit{(1)}~, et %
\textit{(2)}~\textit{ZC} hinc %
\textit{(a)}~secundum %
\textit{(b)}~ponendo~\textit{L}}}
momentum globi totale\indexS{momentum globi totale} esse unitatem,
\edtext{tunc}{\lemma{}\Bfootnote{tunc \textit{erg.}~\textit{L}}} 
secundum
\href{\Scan{LBr_724_16av}}%
{\lbrack16a~v\textsuperscript{o}\rbrack} 
\edtext{hypothesin primae causae applicatae ad planum}{\lemma{hypothesin}\Bfootnote{%
\textit{(1)}~prima causa %
\textit{(a)}~in plano %
\textit{(b)}~applicata ad planum %
\textit{(2)}~primae causae applicatae ad planum~\textit{L}}}
\textit{XC}, qua premitur
\edtext{ipsum}{\lemma{}\Bfootnote{ipsum \textit{erg.}~\textit{L}}}
\edtext{\textit{XC} momento \rule[-1.4ex]{0pt}{5.4ex}$\displaystyle\frac{XC - XN}{XC}$,}{\lemma{\textit{XC}}\Bfootnote{%
\textit{(1)}~momentum %
\textit{(2)}~momento %
\textit{(a)}~ut %
\textit{(b)}~%
\rule[-2.4ex]{0pt}{6ex}%
$\displaystyle\frac{XC - XN}{XC}$,~\textit{L}}}
reliquum momenti totalis \rule[-2.4ex]{0pt}{5.4ex}$\displaystyle\frac{XN}{XC}$ secundum Hypothesin secundae 
\edtext{causae premet planum \textit{ZC}, in quantum}{\lemma{causae}\Bfootnote{%
\textit{(1)}~aget %
\textit{(2)}~premet planum \textit{ZC}. %
\textit{(aa)}~Et contra secundum hypothesin secundae causae in plano %
\textit{(bb)}~, in quantum~\textit{L}}}
enim non
\edtext{sustentatur globus a plano \textit{XC}}{\lemma{sustentatur}\Bfootnote{%
\textit{(1)}~ab uno plano %
\textit{(2)}~globus a plano \textit{XC}~\textit{L}}}
in tantum
\edtext{decliviter descendendo}{\lemma{}\Bfootnote{decliviter descendendo \textit{erg.}~\textit{L}}} 
agit in
\edtext{alterum \textit{ZC.} Et similiter alia methodo secundum}{\lemma{alterum}\Bfootnote{%
\textit{(1)}~, Et %
\textit{(2)}~\textit{ZC}. Et %
\textit{(a)}~con %
\textit{(b)}~similiter \bars{alia methodo \textit{erg.}} secundum~\textit{L}}}
\edtext{hypothesin secundae causae applicatae ad planum}{\lemma{hypothesin}\Bfootnote{%
\textit{(1)}~secundam %
\textit{(2)}~secundae causae %
\textit{(a)}~in plano %
\textit{(b)}~applicatae ad planum~\textit{L}}} 
\textit{XC} quae premitur \textit{XC} momento
\edtext{\rule[-1.4ex]{0pt}{5.4ex}$\displaystyle\frac{ZO}{ZC}$, reliquum}{%
\lemma{$\displaystyle\frac{ZO}{ZC}$,\rule[-2.8ex]{0pt}{6ex}}%
\Bfootnote{%
\textit{(1)}~premetur %
\textit{(2)}~reliquum~\textit{L}}}
momenti totalis\indexS{momentum totale} \rule[-2.8ex]{0pt}{6ex}$\displaystyle\frac{ZC - ZO}{ZC}$
\edtext{secundum hypothesin primae causae}{\lemma{}\Bfootnote{secundum hypothesin primae causae \textit{erg.}~\textit{L}}} 
premet planum\indexS{planum declive} \textit{ZC}, in quantum
\edtext{enim globus decliviter}{\lemma{enim}\Bfootnote{%
\textit{(1)}~non sustentatur globus %
\textit{(2)}~non agit globus in planum %
\textit{(3)}~globus decliviter~\textit{L}}}
descendendo agit in planum \textit{XC},\indexS{globus in planum agens} 
in tantum non sustentatur a plano\indexS{globus a plano non sustentatus}
\edlabel{LBr_724_16a_1a}%
\edtext{}{%
\xxref{LBr_724_16a_1a}{LBr_724_16a_1b}%
\lemma{\textit{ZC}.}%
\Bfootnote{%
\textit{(1)}~Utraque ergo Hypothesis %
\textit{(2)}~Quaelibet ergo methodus %
\textit{(a)}~quae u %
\textit{(b)}~utramque hypothesin jungens~\textit{L}}}%
\textit{ZC}.%
\edlabel{LBr_724_16a_4b}		%
\pend 
\pstart
Quaelibet ergo methodus utramque hypothesin jungens \edlabel{LBr_724_16a_1b}
exhaurit totum momentum globi,\indexS{momentum globi totale} 
sed non possunt ambae ita absolute conjungi, seu non
\edtext{potest unumquodque}{\lemma{potest}\Bfootnote{%
\textit{(1)}~utrumque %
\textit{(2)}~unumquodque~\textit{L}}} 
planum
\edtext{integris}{\lemma{}\Bfootnote{integris \textit{erg.}~\textit{L}}} 
duobus momentis
\edtext{sibi}{\lemma{}\Bfootnote{sibi \textit{erg.}~\textit{L}}} 
assignatis premi, quia ita summa omnium
\edtext{quatuor}{\lemma{}\Bfootnote{quatuor \textit{erg.}~\textit{L}}} 
momentorum partialium\indexS{momentum partiale} aequaretur duplo momento totali.\indexS{momentum totale}
Itaque per \textls{regulam alternativorum},\indexS{regula alternativorum} 
cum nulla sit ratio cur dicatur
\edtext{potius planum \textit{XC} premi momento}{\lemma{potius}\Bfootnote{%
\textit{(1)}~globus %
\textit{(2)}~planum %
\textit{(a)}~moment %
\textit{(b)}~\textit{XC} %
\textit{(aa)}~momenti %
\textit{(bb)}~premi momento~\textit{L}}}
$\displaystyle\frac{XC - XN}{XC}$, quam
\edtext{momento\indexS{momentum}}{\lemma{}\Bfootnote{momento \textit{erg.}~\textit{L}}} 
$\displaystyle\frac{ZO}{ZC}$ seu 
\edtext{$\displaystyle\frac{ZO}{XC}$, utrumque}{%
\lemma{$\displaystyle\frac{ZO}{XC}$,\rule[-2.4ex]{0pt}{6.0ex}}%
\Bfootnote{%
\textit{(1)}~utraque %
\textit{(2)}~utrumque~\textit{L}}}
autem in rationes venire
\edtext{debeat, ideo dimidio amborum premetur planum \textit{XC}, seu}{\lemma{debeat}\Bfootnote{%
\textit{(1)}~premetur \bars{mom \textit{erg.~u.~gestr.}} dimidio amborum, seu %
\textit{(2)}~ideo dimidio amborum premetur planum \textit{XC}, seu~\textit{L}}}
momento
\edtext{\rule[-1.2ex]{0pt}{5.4ex}$\displaystyle\frac{XC - XN + ZO}{2. XC}$. Et similiter}{%
\lemma{$\displaystyle\frac{XC - XN + ZO}{2. XC}$.\rule[-2.7ex]{0pt}{6ex}}%
\Bfootnote{%
\textit{(1)}~Et similiter planum \textit{ZC} momento $XC + XN$ %
\textit{(2)}~Et similiter~\textit{L}}}
planum \textit{ZC} premetur momento \rule[-2.8ex]{0pt}{6.4ex}$\displaystyle\frac{XC - ZO + XN}{2. XC}$
quorum amborum momentorum\indexS{momentum partiale} summa facit momentum totale\indexS{momentum totale} seu unitatem.
\edtext{%
Sit \textit{XC}, 5; \textit{XN}, 3;
\edtext{\textit{ZO}, 4. Momentum}{\lemma{\textit{ZO}, 4.}\Bfootnote{%
\textit{(1)}~fiet %
\textit{(2)}~Momentum~\textit{L}}} 
quo premitur planum \textit{XC} erit \rule[-2.4ex]{0pt}{6ex}$\displaystyle\frac{3}{5}$;
et momentum quo premitur planum \textit{ZC} erit
\edlabel{LBr_724_16a_2a}%
\edtext{}{%
\xxref{LBr_724_16a_2a}{LBr_724_16a_2b}%
\lemma{$\displaystyle\frac{2}{5}$.}%
\Bfootnote{\bars{itaque mom \textit{gestr.}}
\href{\Scan{LBr_724_16ar}}%
{\lbrack16a~r\textsuperscript{o}\rbrack} NB.~\textit{L}}}%
\rule[-2.4ex]{0pt}{6ex}$\displaystyle\frac{2}{5}$.}{
\lemma{}%
\applabel{FnA:LBr_724_16a_Nebenrechnung}%
\Afootnote{%
\textit{Hierzu, am Rand von 
\href{\Scan{LBr_724_16v}}%
{Bl.~16~v\textsuperscript{o}}:} $\displaystyle\frac{4}{5}+\displaystyle\frac{3}{5}$}}
\pend
\pstart
\href{\Scan{LBr_724_16ar}}%
{\lbrack16a~r\textsuperscript{o}\rbrack} 
NB.\edlabel{LBr_724_16a_2b} Cum sit \textit{B} ad \textit{C} ut \textit{AB} ad \textit{AC}\lbrack,\rbrack\
\edtext{hinc tota vis\indexS{vis globum deprimens}}{%
\lemma{hinc}%
\Bfootnote{%
\bars{de \textit{streicht Hrsg.}} tota vis~\textit{L}}}
globum \textit{C} deprimens sustinebitur vi globum \textit{B} deprimente,\indexS{vis globum deprimens} 
quae est ipsius \rule[-1.4ex]{0pt}{5.4ex}$\displaystyle\frac{AB}{AC}$, et elaterio\indexS{elaterium plani} plani in \textit{C} 
quod proinde patitur vim ut \rule[-2.8ex]{0pt}{6ex}$\displaystyle\frac{AC - AB}{AC}$\lbrack,\rbrack\
et tanta vi censendus est globus\indexS{globus in planum agens} \textit{C} perpendiculariter agere in planum \textit{AC}.
\pend
\begin{edfigure}[0.5em]%
\includegraphics[scale=1]{%
gesamttex/edit_VIII-4/images/LBr_724_016a_d2_016ar.pdf}
\edcaption{\lbrack\textit{Fig.~2}\rbrack}
\end{edfigure}
\pstart
Quae actio obliqua\indexS{actio obliqua potentiae mortuae} potentiae mortuae,\indexS{potentia mortua} 
longe abest ab actione obliqua potentiae\indexS{actio obliqua potentiae vivae} vivae\indexS{potentia viva} seu impetus\indexS{impetus} quaesiti.
\pend
